\documentclass[11pt,a4paper]{article}
\usepackage{fontspec,amsmath,amssymb,amsthm,microtype}
\setmainfont{lmroman10-regular.otf}[BoldFont=lmroman10-bold.otf,ItalicFont=lmroman10-italic.otf,BoldItalicFont=lmroman10-bolditalic.otf]
\setmonofont{lmmono10-regular.otf}[Scale=MatchLowercase]
\usepackage[margin=25mm]{geometry}
\usepackage[unicode,hidelinks]{hyperref}
\setlength{\emergencystretch}{3em}
\setlength{\parskip}{0.4em}
\setlength{\parindent}{0pt}
\allowdisplaybreaks[2]
\raggedbottom
\title{\texorpdfstring{Conditional LSI theorem for the actual ULA laws}{Conditional LSI theorem for the actual ULA laws}}
\author{LangevinLean mathematical exposition}
\date{}
\begin{document}\maketitle
\noindent Source module: \texttt{Unadjusted\allowbreak{}Langevin/\allowbreak{}LSI/\allowbreak{}Theorem.\allowbreak{}lean}. The accompanying HTML pairs this mathematical exposition with the exact Lean statements and records the scope of the formalization.


The actual iterated Gaussian ULA laws satisfy \(H_k\le e^{-\alpha hk}H_0+8hdL^2/\alpha\), once the displayed interpolation identification, finite KL, entropy/\allowbreak{}Fisher dissipation, LSI, and moment hypotheses are supplied. This endpoint does not construct all those analytic data from LSI and smoothness alone.



The full actual Gaussian ULA endpoint in UnconditionalTheorem.lean derives all necessary analytic estimates from smoothness, LSI, and finite initial KL. This module retains a reusable interface for explicitly specified interpolation data.


\section{\texorpdfstring{The target estimate and the conditional Lean statement}{The target estimate and the conditional Lean statement}}


For every index, derive the one-step recurrence from the preceding interpolation theorem. The numerical conditions \(\alpha\le L\), \(h\le\alpha/(4L^2)\) imply \(\alpha h\le1/4\). Apply \texttt{discrete\_\allowbreak{}entropy\_\allowbreak{}recurrence\_\allowbreak{}bound} and rewrite the initial value with \texttt{ula\allowbreak{}Entropy\_\allowbreak{}zero}. This gives \(H_k\le e^{-\alpha hk}\operatorname{KL}(\mu\mid\pi_U)_{\mathbb R}+8hdL^2/\alpha\), conditional on the displayed analytic path hypotheses.



{\raggedright\noindent Lean declarations: \texttt{ula\_\allowbreak{}kl\_\allowbreak{}bound\_\allowbreak{}of\_\allowbreak{}interpolation}.\par}

\section{\texorpdfstring{Why finite KL and Fisher information are separate obligations}{Why finite KL and Fisher information are separate obligations}}


At time zero the supplied relative-density law identity gives \(r_{k,0}\nu=\rho_k\), because \texttt{interpolation\allowbreak{}Law\_\allowbreak{}zero} returns its input law. Substitute this equality into finite KL of the time-zero density path. It yields finite actual KL of the algorithm law at every index.



{\raggedright\noindent Lean declarations: \texttt{ula\_\allowbreak{}kl\_\allowbreak{}finite\_\allowbreak{}of\_\allowbreak{}interpolation}.\par}

\section{\texorpdfstring{Combine production, LSI, and the initial gradient moment}{Combine production, LSI, and the initial gradient moment}}


Define \(G_k=\int\|\nabla U(x)\|^2d\rho_k(x)\) for the actual ULA law. Since Lean's ordinary integral is totalized, later hypotheses explicitly prove integrability of the squared gradient before treating \(G_k\) as an energy. The definition itself contains no entropy-moment estimate.



{\raggedright\noindent Lean declarations: \texttt{ula\allowbreak{}Gradient\allowbreak{}Moment}.\par}

\section{\texorpdfstring{The direct endpoint constructs the recurrence from the path hypotheses}{The direct endpoint constructs the recurrence from the path hypotheses}}


Identify the entropy path's endpoints using \texttt{interpolation\allowbreak{}Law\_\allowbreak{}zero} and \texttt{interpolation\allowbreak{}Law\_\allowbreak{}endpoint}: its time-zero entropy is \(H_k\), and its time-\(h\) entropy is \(H_{k+1}\). The law identity and probability preservation prove that every intermediate relative density is normalized.



Apply the real LSI inequality using the explicit positive smooth density and finite KL/\allowbreak{}Fisher guards. The scalar \texttt{one\_\allowbreak{}step\_\allowbreak{}entropy\_\allowbreak{}of\_\allowbreak{}interpolation} then combines it with the supplied entropy derivative, production, and gradient-moment estimates. Rewriting the two endpoints yields \(H_{k+1}\le e^{-\alpha h}H_k+6h^2dL^2\). Those analytic interpolation premises remain explicit in this theorem.



{\raggedright\noindent Lean declarations: \texttt{ula\_\allowbreak{}kl\_\allowbreak{}one\_\allowbreak{}step\_\allowbreak{}of\_\allowbreak{}interpolation}.\par}

\end{document}
